% Part: incompleteness
% Chapter: representability-in-q
% Section: representing-relations

\documentclass[../../../include/open-logic-section]{subfiles}

\begin{document}

% SEGMENT OLP-0298-S01

\olfileid{inc}{req}{rel}

\olsection{உறவுகளைப் பிரதிநிதித்துவப்படுத்தல்}

ஓர் \emph{உறவு} பிரதிநிதித்துவப்படுத்தத்தக்கது என்பதன் பொருளை இப்போது
கூறுவோம்.

\begin{defn}
\ollabel{defn:representing-relations} இயல் எண்கள் மீதான ஓர் உறவு
$R(x_0,\dots,x_k)$ பின்வரும் நிபந்தனையை நிறைவுசெய்தால்,
{\em $\Th{Q}$ இல் பிரதிநிதித்துவப்படுத்தத்தக்கது} எனப்படுகிறது:
$!A_R(x_0,\dots,x_k)$ என்ற வாய்பாடு ஒன்று இருந்து,
$R(n_0,\dots,n_k)$ மெய்யாக இருக்கும்போதெல்லாம் $\Th{Q}$ ஆனது
$!A_R(\num{n_0},\dots,\num{n_k})$ ஐ நிறுவுகிறது; மேலும்
$R(n_0,\dots,n_k)$ பொய்யாக இருக்கும்போதெல்லாம் $\Th{Q}$ ஆனது
$\lnot !A_R(\num{n_0}, \dots, \num{n_k})$ ஐ நிறுவுகிறது.
\end{defn}

\begin{thm}
\ollabel{thm:representing-rels} ஓர் உறவு~$\Th{Q}$ இல்
பிரதிநிதித்துவப்படுத்தத்தக்கது எனில், எனில் மட்டுமே அது கணிக்கத்தக்கது.
\end{thm}

\begin{proof}
முன்திசைக்கு, $R(x_0,\dots,x_k)$ என்ற உறவு
$!A_R(x_0,\dots,x_k)$ என்ற வாய்பாட்டால்
பிரதிநிதித்துவப்படுத்தப்படுகிறது என்று கொள்க. $R$ ஐக் கணிப்பதற்கான
ஒரு படிமுறை இதோ: $n_0$, \dots,~$n_k$ ஆகியவற்றை உள்ளீடாகக் கொண்டு,
$!A_R(\num{n_0}, \dots, \num{n_k})$ இன் ஓர் நிறுவலையும்
$\lnot !A_R(\num{n_0}, \dots, \num{n_k})$ இன் ஓர் நிறுவலையும் ஒரே நேரத்தில்
தேடுக. நமது எடுகோளின்படி இத்தேடல் அவற்றுள் ஒன்றைத் தவறாமல் கண்டடையும்;
முதலாவதைக் கண்டால் “ஆம்” என்றும், இல்லையெனில் “இல்லை” என்றும்
விடையளிக்க.

மறுதிசையில், $R(x_0, \dots, x_k)$ கணிக்கத்தக்கது என்று கொள்க.
வரையறையின்படி, இதன் பொருள் $\Char{R}(x_0, \dots, x_k)$ என்ற சார்பு
கணிக்கத்தக்கது என்பதாகும். \olref[int]{thm:representable-iff-comp}
இன்படி, $\Char{R}$ ஆனது ஒரு வாய்பாட்டால் பிரதிநிதித்துவப்படுத்தப்படுகிறது;
அதை $!A_{\Char{R}}(x_0, \dots, x_k, y)$ என்க. $!A_R(x_0, \dots, x_k)$
என்பது $!A_{\Char{R}}(x_0, \dots, x_k, \num{1})$ என்ற வாய்பாடாக
இருக்கட்டும். அப்போது எந்த $n_0$, \dots,~$n_k$ க்கும்,
$R(n_0, \dots, n_k)$ மெய்யாக இருந்தால்
$\Char{R}(n_0, \dots, n_k) = 1$. இந்நேர்வில் $\Th{Q}$ ஆனது
$!A_{\Char{R}}(\num{n_0}, \dots, \num{n_k}, \num{1})$ ஐ நிறுவுகிறது;
எனவே $\Th{Q}$ ஆனது $!A_R(\num{n_0}, \dots, \num{n_k})$ ஐ நிறுவுகிறது.
மறுபுறம், $R(n_0, \dots, n_k)$ பொய்யாக இருந்தால்
$\Char{R}(n_0, \dots, n_k) = 0$. இதனால் $\Th{Q}$ பின்வருவதை நிறுவுகிறது:
\[
\lforall[y][(!A_{\Char{R}}(\num{n_0}, \dots, \num{n_k}, y) \lif y =
  \num{0})].
\]
$\Th{Q}$ ஆனது $\eq/[\num{0}][\num{1}]$ ஐ நிறுவுவதால்,
$\Th{Q}$ ஆனது
$\lnot !A_{\Char{R}}(\num{n_0}, \dots, \num{n_k}, \num{1})$ ஐ
நிறுவுகிறது; எனவே அது
$\lnot !A_R(\num{n_0}, \dots, \num{n_k})$ ஐ நிறுவுகிறது.
\end{proof}

\begin{prob}
$R$ ஆனது~$\Th{Q}$ இல் பிரதிநிதித்துவப்படுத்தத்தக்கது எனில்,
$\Char{R}$ உம் அவ்வாறே என்பதை நிறுவுக.
\end{prob}

\end{document}
